\documentclass[12pt]{article}
\usepackage[T1]{fontenc}
\usepackage[utf8]{inputenc}
\usepackage{lmodern}
\usepackage{textcomp}
\usepackage{enumitem}
\usepackage{geometry}
\usepackage{graphicx}
\usepackage{amsmath, amssymb}
\geometry{margin=1in}
\begin{document}
\begin{center}
Chemistry - Graduate Level Assessment 2 \
Total Marks: 40 \
Answer all questions.
\end{center}

\section*{Multiple Choice (10 marks)}
\begin{enumerate}[label=\arabic*.]
\item Elements with partially filled 4 f or 5 f orbitals include all of the following EXCEPT
\begin{enumerate}[label=(\alph*)]
    \item Eu
    \item Gd
    \item Am
    \item Cu
\end{enumerate}
\item Assuming complete neutralization, calculate the number of milliliters of 0.025 M H\_3 PO\_4 required to neutralize 25 ml of 0.030 M Ca(OH)\_2.
\begin{enumerate}[label=(\alph*)]
    \item 10 ml
    \item 30 ml
    \item 55 ml
    \item 40 ml
    \item 25 ml
\end{enumerate}
\item A wine has an acetic acid (CH\_3 COOH, 60 g/formula weight) content of 0.66\% by weight. If the density of the wine is 1.11 g/ml, what is the formality of the acid?
\begin{enumerate}[label=(\alph*)]
    \item 1.6 \ensuremath{\times} 10^-3 F
    \item 1.5 \ensuremath{\times} 10^-3 F
    \item 1.4 \ensuremath{\times} 10^-3 F
    \item 0.8 \ensuremath{\times} 10^-3 F
    \item 0.9 \ensuremath{\times} 10^-3 F
\end{enumerate}
\item Under standard temperature and pressure conditions, compare the relative rates at which inert gases,Ar, He, and Kr diffuse through a common orifice.
\begin{enumerate}[label=(\alph*)]
    \item .1002 : .3002 : .4002
    \item .3582 : .4582 : .0582
    \item .2582 : .4998 : .3092
    \item .1582 : .6008 : .2092
    \item .1582 : .4998 : .1092
\end{enumerate}
\item Would you expect He\_2^\ensuremath{\oplus} to be more stable than He\_2? Than He?
\begin{enumerate}[label=(\alph*)]
    \item He\_2^\ensuremath{\oplus} is the most stable among He, He\_2, and He\_2^\ensuremath{\oplus}
    \item He\_2^\ensuremath{\oplus} is more stable than He\_2 but less stable than He
    \item He, He\_2, and He\_2^\ensuremath{\oplus} have the same stability
    \item He\_2^\ensuremath{\oplus} is the least stable among He, He\_2, and He\_2^\ensuremath{\oplus}
    \item He and He\_2 have the same stability, both more stable than He\_2^\ensuremath{\oplus}
\end{enumerate}
\end{enumerate}

\section*{True / False (10 marks)}
\textit{Answer True or False}
\begin{enumerate}[label=\arabic*.]
\setcounter{enumi}{5}
\item True or False: The ionization energy of singly ionized helium calculated using the Bohr theory is 52.307 eV.
\item True or False: At$_{12}$.19 K, a pure sample of an ideal gas can exhibit a pressure of 1 atm and a concentration of 1 mole liter^-1.
\item True or False: NiSO 3 is expected to have low solubility in water due to its ionic nature and the limited solubility of sulfites in aqueous solutions.
\item True or False: The missing mass in ^19\_9 F is 0.00549 amu, which suggests a significant mass loss due to nuclear reactions.
\item True or False: The molality of a solution containing 49 g of H$_{2}$SO$_{4}$ dissolved in 250 g of water is 2.0 moles/kg.
\end{enumerate}

\section*{Long Form Response (20 marks)}
\textit{Answer each question with a clear and complete explanation.}
\begin{enumerate}[label=\arabic*.]
\setcounter{enumi}{10}
\item Discuss the process of diluting a 90\% ammonia solution to achieve a 60\% concentration. How much water is required, and what calculations support your answer?
\item Discuss the thermodynamic principles governing the expansion of one cubic foot of air when heated from 0^\ensuremath{\circC} to 546,000^\ensuremath{\circC}, and calculate the resulting volume.
\end{enumerate}

\vfill
\noindent\textit{End of Paper}
\end{document}